% =============================================================================
% Chapter 4: Implementation
% =============================================================================
% Required packages (add to your main preamble):
%   \usepackage{listings}
%   \usepackage{booktabs}
%   \usepackage{tikz}
%   \usetikzlibrary{shapes, arrows.meta, positioning}
%   \usepackage{dirtree}  % optional, for directory trees
% =============================================================================

\chapter{Implementation}
\label{ch:implementation}

This chapter describes the implementation details of the broker and agent components, the technology choices, project structure, and the integration with Liqo for automatic cluster peering.


% =============================================================================
\section{Technology Stack}
\label{sec:tech-stack}

Table~\ref{tab:tech-stack} summarizes the key technologies used in the implementation and the rationale for each choice.

\begin{table}[htbp]
\centering
\caption{Technology stack}
\label{tab:tech-stack}
\begin{tabular}{@{}llp{6.5cm}@{}}
\toprule
\textbf{Technology} & \textbf{Version} & \textbf{Rationale} \\
\midrule
Go & 1.24 & Primary language of the Kubernetes ecosystem; strong concurrency, static typing, fast compilation \\
controller-runtime & 0.22 & Standard framework for building Kubernetes operators; provides reconciliation loop, caching, leader election \\
Kubebuilder & 4.x & Scaffolding tool for CRD definitions and controller boilerplate \\
cert-manager & 1.x & Kubernetes-native certificate lifecycle management; automates CA creation and certificate rotation \\
Kind & 0.20+ & Kubernetes-in-Docker for local multi-cluster testing \\
Liqo & latest & Multi-cluster peering via virtual nodes and WireGuard tunnels \\
Docker & 24+ & Container runtime for Kind clusters \\
\bottomrule
\end{tabular}
\end{table}


% =============================================================================
\section{Project Structure}
\label{sec:project-structure}

The project is organized as a monorepo containing two independent Go modules: the broker and the agent. In a production deployment these would be separate repositories, but they are combined here to simplify testing and evaluation.

\subsection{Broker Structure}

\begin{verbatim}
liqo-resource-broker/
├── api/v1alpha1/              # CRD type definitions
│   ├── clusteradvertisement_types.go
│   └── reservation_types.go
├── cmd/main.go                # Entry point, flag parsing, server startup
├── internal/
│   ├── api/                   # HTTP layer
│   │   ├── server.go          # TLS server setup and route registration
│   │   ├── handlers/          # POST/GET handlers for each endpoint
│   │   └── middleware/        # mTLS authentication, logging
│   ├── broker/
│   │   └── decision.go        # Decision engine (filter, score, select)
│   ├── controller/
│   │   └── reservation_controller.go  # Reconciler for Reservation CRD
│   └── resource/
│       └── availability.go    # Available = Allocatable - Allocated - Reserved
└── config/
    ├── crd/                   # Generated CRD YAML manifests
    └── certmanager/           # Certificate and Issuer definitions
\end{verbatim}

\subsection{Agent Structure}

\begin{verbatim}
liqo-resource-agent/
├── api/v1alpha1/              # CRD type definitions
│   ├── advertisement_types.go
│   ├── reservationinstruction_types.go
│   ├── providerinstruction_types.go
│   └── resourcerequest_types.go    # User-facing reservation trigger
├── cmd/main.go                # Entry point, flag parsing, controller setup
├── internal/
│   ├── controller/
│   │   ├── advertisement_controller.go       # Reconciles local Advertisement
│   │   ├── resourcerequest_controller.go     # Synchronous reservation flow
│   │   ├── reservationinstruction_controller.go  # Handles instructions + Liqo
│   │   ├── providerinstruction_controller.go # Handles provider instructions
│   │   └── instruction_poller.go   # Polls GET /instructions every 5s
│   ├── metrics/
│   │   └── collector.go       # Node/pod resource collection
│   ├── publisher/
│   │   └── broker_client.go   # Legacy Kubernetes transport
│   └── transport/
│       ├── interface.go       # BrokerCommunicator interface
│       └── http/
│           └── client.go      # mTLS HTTP client with retry logic
└── config/
    └── crd/                   # Generated CRD YAML manifests
\end{verbatim}


% =============================================================================
\section{Broker Implementation}
\label{sec:broker-impl}

% -----------------------------------------------------------------------------
\subsection{API Server}
\label{subsec:api-server}

The broker's HTTP server is implemented using Go's standard \texttt{net/http} package with TLS configuration for mutual authentication. The server initialization proceeds as follows:

\begin{enumerate}
    \item Load the server certificate and private key from the filesystem.
    \item Load the Certificate Authority (CA) certificate into a \texttt{x509.CertPool}.
    \item Configure \texttt{tls.Config} with \texttt{ClientAuth: tls.RequireAndVerifyClientCert} and the CA pool as \texttt{ClientCAs}.
    \item Register API routes on an \texttt{http.ServeMux}.
    \item Apply the middleware chain: authentication first, then logging.
    \item Start the HTTPS listener.
\end{enumerate}

The middleware chain is applied in reverse order using a \texttt{Chain} function, ensuring that authentication executes before the request reaches any handler.

% -----------------------------------------------------------------------------
\subsection{Authentication Middleware}
\label{subsec:auth-middleware}

The authentication middleware extracts the cluster identity from the client certificate's Common Name (CN):

\begin{lstlisting}[language=Go, caption={Client certificate validation middleware (simplified)}, label={lst:auth-middleware}]
func ValidateClientCertificate(next http.Handler) http.Handler {
    return http.HandlerFunc(func(w http.ResponseWriter, r *http.Request) {
        if r.URL.Path == "/healthz" {
            next.ServeHTTP(w, r)
            return
        }
        if r.TLS == nil || len(r.TLS.PeerCertificates) == 0 {
            http.Error(w, "client certificate required", 401)
            return
        }
        clusterID := r.TLS.PeerCertificates[0].Subject.CommonName
        ctx := context.WithValue(r.Context(), ClusterIDKey, clusterID)
        next.ServeHTTP(w, r.WithContext(ctx))
    })
}
\end{lstlisting}

Downstream handlers retrieve the cluster ID from the request context. The \texttt{PostAdvertisement} handler additionally validates that the \texttt{clusterID} in the request body matches the certificate CN, preventing a compromised agent from publishing data on behalf of another cluster.

% -----------------------------------------------------------------------------
\subsection{Advertisement Handler}
\label{subsec:adv-handler}

The \texttt{PostAdvertisement} handler receives a JSON-encoded advertisement from an agent and creates or updates the corresponding \texttt{ClusterAdvertisement} CRD in the broker cluster. A critical implementation detail is the \textbf{preservation of the Reserved field}:

\begin{enumerate}
    \item The handler fetches the existing \texttt{ClusterAdvertisement} (if any).
    \item If found, it copies the existing \texttt{Resources.Reserved} value before overwriting the resource fields with the new data from the agent.
    \item The updated resource (with preserved \texttt{Reserved}) is written back.
\end{enumerate}

This ensures that the broker's resource locking decisions are not lost when an agent publishes a routine advertisement update. Without this preservation, a race condition would occur: the broker locks resources (sets \texttt{Reserved}), and then the agent's next publish cycle overwrites \texttt{Reserved} with zero, effectively unlocking resources prematurely.

The advertisement name in the broker follows the convention \texttt{<clusterID>-adv} (e.g., \texttt{agent-cluster-1-adv}).

% -----------------------------------------------------------------------------
\subsection{Reservation Handler}
\label{subsec:reservation-handler}

The \texttt{PostReservation} handler processes synchronous reservation requests from agents. When an agent sends a \texttt{POST /api/v1/reservations} with the requested CPU, memory, and priority, the handler:

\begin{enumerate}
    \item Validates the request and extracts the requester's cluster ID from the mTLS certificate CN.
    \item Creates a \texttt{Reservation} CRD in the broker cluster with phase \textit{Pending}.
    \item Invokes the decision engine's \texttt{SelectBestCluster} method to find the optimal provider.
    \item If a provider is found, locks resources using \texttt{RetryOnConflict} and transitions the reservation to \textit{Reserved}.
    \item Returns the reservation result (including the selected provider and reservation name) in the HTTP response.
\end{enumerate}

This inline processing ensures that the requester receives its instruction in a single HTTP round trip, eliminating the need for a separate polling mechanism on the requester side.

The \texttt{GetInstructions} handler serves the \texttt{GET /api/v1/instructions} endpoint. It lists all \textit{Reserved}-phase reservations where the calling cluster (identified by mTLS certificate CN) is the target provider, and returns them as a JSON array.

% -----------------------------------------------------------------------------
\subsection{Reservation Controller}
\label{subsec:reservation-controller}

The \texttt{ReservationReconciler} is a Kubernetes controller that watches \texttt{Reservation} resources and drives them through their lifecycle. It implements the standard controller-runtime reconciliation pattern.

\subsubsection{Reconciliation Logic}

The reconciler handles each phase as follows:

\begin{itemize}
    \item \textbf{Pending:} Invokes the decision engine's \texttt{SelectBestCluster} method. If a provider is found, the reconciler locks resources using \texttt{retry.RetryOnConflict} to handle concurrent updates, then transitions the reservation to \textit{Reserved}. If no suitable cluster exists, the phase becomes \textit{Failed}.

    \item \textbf{Reserved:} Monitors for agent acknowledgment (via a status condition). If the requester agent confirms activation, the phase transitions to \textit{Active}. If the reservation expires, resources are released and the phase becomes \textit{Released}.

    \item \textbf{Active:} Monitors for explicit release signals or expiration. When either occurs, the controller calls \texttt{releaseResources} to decrement the provider's \texttt{Reserved} field and transitions to \textit{Released}.

    \item \textbf{Failed / Released:} Terminal states with no further action.
\end{itemize}

A \texttt{Finalizer} ensures that resources are properly released even if the reservation is deleted directly:

\begin{lstlisting}[language=Go, caption={Finalizer-based resource cleanup (simplified)}, label={lst:finalizer}]
if !reservation.DeletionTimestamp.IsZero() {
    if err := r.releaseResources(ctx, reservation); err != nil {
        return ctrl.Result{}, err
    }
    controllerutil.RemoveFinalizer(reservation, ReservationFinalizer)
    return ctrl.Result{}, r.Update(ctx, reservation)
}
\end{lstlisting}

\subsubsection{Resource Locking}

The resource locking procedure uses Kubernetes' optimistic concurrency control via \texttt{retry.RetryOnConflict}:

\begin{enumerate}
    \item Fetch the latest version of the target \texttt{ClusterAdvertisement}.
    \item Verify sufficient resources are still available (another reservation may have consumed them).
    \item Add the requested quantities to the \texttt{Reserved} field.
    \item Attempt to update the resource. If a conflict occurs (another concurrent update changed the resource version), retry from step 1.
\end{enumerate}

This approach guarantees consistency without requiring external locking mechanisms.

% -----------------------------------------------------------------------------
\subsection{Decision Engine}
\label{subsec:decision-engine}

The decision engine implements the algorithm described in Section~\ref{sec:decision-algorithm}. The key implementation function is \texttt{SelectBestCluster}, which accepts the requester's cluster ID, requested CPU, requested memory, and priority, then returns the best matching \texttt{ClusterAdvertisement} or an error.

The base score for display purposes (stored in \texttt{ClusterAdvertisement.Status.Score}) is computed on a 0--100 scale:

\begin{equation}
\text{BaseScore} = \frac{\text{Available}_{\text{cpu}}}{\text{Allocatable}_{\text{cpu}}} \times 50 + \frac{\text{Available}_{\text{mem}}}{\text{Allocatable}_{\text{mem}}} \times 50
\end{equation}

This provides a quick visual indicator of cluster health: a score of 100 means the cluster is fully available, while 0 means it is fully utilized.


% =============================================================================
\section{Agent Implementation}
\label{sec:agent-impl}

% -----------------------------------------------------------------------------
\subsection{Resource Collector}
\label{subsec:resource-collector}

The \texttt{Collector} struct is responsible for gathering accurate resource metrics from the local cluster. It implements the \texttt{CollectClusterResources} method, which performs the following steps:

\begin{enumerate}
    \item \textbf{Node enumeration:} Lists all nodes and filters for ready nodes using the \texttt{NodeReady} condition.

    \item \textbf{Capacity and allocatable:} For each ready node, sums the \texttt{Capacity} (total hardware) and \texttt{Allocatable} (capacity minus system reservations) values for CPU, memory, and GPU.

    \item \textbf{Allocated resources:} Lists all pods and sums their resource \textit{requests} (not limits). For init containers, the maximum across all init containers is used rather than the sum, since init containers run sequentially. Pod overhead is also included.

    \item \textbf{Reserved resources:} Lists all \texttt{ProviderInstruction} CRDs that are enforced and not expired, summing their requested CPU and memory. This accounts for resources locked by the broker for other clusters.

    \item \textbf{Available computation:} Computes \texttt{Available = Allocatable - Allocated - Reserved}, clamping negative values to zero.
\end{enumerate}

The collector runs within the \texttt{AdvertisementReconciler}'s reconciliation loop, triggered periodically by the configurable \texttt{--advertisement-requeue-interval} (default: 30 seconds).

% -----------------------------------------------------------------------------
\subsection{HTTP Transport}
\label{subsec:http-client}

The \texttt{HTTPCommunicator} implements the \texttt{BrokerCommunicator} interface, which defines the protocol-agnostic contract for agent-broker communication:

\begin{lstlisting}[language=Go, caption={BrokerCommunicator interface}, label={lst:communicator}]
type BrokerCommunicator interface {
    PublishAdvertisement(ctx context.Context,
        adv *dto.AdvertisementDTO) error
    RequestReservation(ctx context.Context,
        req *dto.ReservationRequestDTO) (*dto.ReservationDTO, error)
    FetchInstructions(ctx context.Context) (
        []*dto.ReservationDTO, error)
    Ping(ctx context.Context) error
    Close() error
}
\end{lstlisting}

The \texttt{RequestReservation} method sends a synchronous \texttt{POST /api/v1/reservations} and returns the broker's decision in the response. The \texttt{FetchInstructions} method performs a lightweight \texttt{GET /api/v1/instructions} to discover provider instructions. This interface enables adding new transport protocols (MQTT, gRPC) without modifying the controllers.

\subsubsection{mTLS Client Configuration}

The HTTP communicator creates a TLS-configured \texttt{http.Client} with:
\begin{itemize}
    \item Client certificate and key loaded from the filesystem.
    \item The broker's CA certificate as the trusted root.
    \item TLS 1.2 minimum version.
    \item Connection pooling: 10 max idle connections, 90-second idle timeout.
    \item Request timeout: 30 seconds.
\end{itemize}

\subsubsection{Reserved Field Preservation}

Before publishing a new advertisement, the HTTP communicator first fetches the current advertisement from the broker via \texttt{GET /api/v1/advertisements/\{clusterID\}}. If the response contains a non-nil \texttt{Reserved} field, the communicator copies this value into the outgoing advertisement. This two-step process ensures that the agent never accidentally overwrites broker-side resource locks.

\subsubsection{Retry with Exponential Backoff}

The \texttt{doWithRetry} method implements retry logic for transient failures:
\begin{itemize}
    \item Maximum retries: 3 attempts.
    \item Initial backoff: 1 second, doubling each attempt, capped at 16 seconds.
    \item Only server errors (HTTP 5xx) and network errors trigger retries.
    \item Client errors (4xx) are not retried.
\end{itemize}

% -----------------------------------------------------------------------------
\subsection{Synchronous Reservation Flow}
\label{subsec:sync-reservation}

The \texttt{ResourceRequestReconciler} watches for \texttt{ResourceRequest} CRDs created by users. When a new request appears, the controller:

\begin{enumerate}
    \item Calls \texttt{RequestReservation} on the \texttt{BrokerCommunicator}, which sends a synchronous \texttt{POST /api/v1/reservations} to the broker.
    \item The broker processes the request inline: running the decision engine, locking resources via \texttt{RetryOnConflict}, and creating the \texttt{Reservation} CRD.
    \item The broker returns the reservation result (including the selected provider and reservation name) in the HTTP response.
    \item If the reservation succeeded, the controller creates a local \texttt{ReservationInstruction} CRD and updates the \texttt{ResourceRequest} status to \textit{Reserved} with the target cluster ID.
    \item If the reservation failed (no suitable provider), the status is set to \textit{Failed} with a descriptive message.
\end{enumerate}

This synchronous design eliminates polling latency for the requester: the instruction is available immediately after the HTTP round trip, typically under 1 second.

% -----------------------------------------------------------------------------
\subsection{Instruction Polling}
\label{subsec:instruction-polling}

The \texttt{InstructionPoller} runs as a background goroutine that discovers provider instructions for this cluster. It queries \texttt{GET /api/v1/instructions} every 5 seconds (configurable via \texttt{--instruction-poll-interval}). The broker returns all \textit{Reserved}-phase reservations where this cluster is the target provider.

For each returned reservation, the poller creates a local \texttt{ProviderInstruction} CRD if one does not already exist. This lightweight polling mechanism ensures that provider clusters discover new instructions within at most 5 seconds of the broker's decision.

% -----------------------------------------------------------------------------
\subsection{Instruction Controllers}
\label{subsec:instruction-controllers}

Two controllers handle the locally created instruction CRDs:

\begin{itemize}
    \item \textbf{ReservationInstructionReconciler:} Processes instructions on the requester side. When a new instruction arrives (with \texttt{delivered = false}), the controller triggers Liqo peering (described in Section~\ref{sec:liqo-impl}) and marks the instruction as delivered. It also monitors for expiration and marks expired instructions accordingly.

    \item \textbf{ProviderInstructionReconciler:} Processes instructions on the provider side. It logs the reservation details and marks the instruction as enforced. The enforced instructions are included in the agent's resource calculation (\texttt{Reserved} component) via the \texttt{Collector}.
\end{itemize}


% =============================================================================
\section{Liqo Integration}
\label{sec:liqo-impl}

The Liqo integration bridges the gap between the broker's resource \emph{decision} and the actual cluster \emph{interconnection}. When the requester agent receives a \texttt{ReservationInstruction}, it automatically establishes peering with the provider.

\subsection{Configuration}

The agent accepts a \texttt{--kubeconfigs-dir} flag specifying a directory containing kubeconfig files for all clusters in the federation. Each file is named \texttt{<clusterID>.kubeconfig}. This directory is populated during cluster setup (e.g., by exporting kubeconfigs from Kind clusters).

If \texttt{--kubeconfigs-dir} is not set, Liqo peering is skipped and the instruction is marked as delivered immediately. This allows the system to operate without Liqo for testing or demonstration purposes.

\subsection{Peering Execution}

The \texttt{executeLiqoPeering} method in the \texttt{ReservationInstructionReconciler} performs the following steps:

\begin{enumerate}
    \item Constructs the local kubeconfig path: \texttt{<dir>/<ownClusterID>.kubeconfig}.
    \item Constructs the remote kubeconfig path: \texttt{<dir>/<targetClusterID>.kubeconfig}.
    \item Verifies both files exist and that \texttt{liqoctl} is available in the system PATH.
    \item Executes \texttt{liqoctl peer} with a 5-minute timeout:
\end{enumerate}

\begin{lstlisting}[language=bash, caption={Liqo peering command}, label={lst:liqo-peer}]
liqoctl peer \
    --kubeconfig <local-kubeconfig> \
    --remote-kubeconfig <remote-kubeconfig> \
    --gw-server-service-type NodePort
\end{lstlisting}

The \texttt{--gw-server-service-type NodePort} flag is required for Kind clusters, which do not support the \texttt{LoadBalancer} service type. In a cloud deployment, this flag would be omitted or changed to \texttt{LoadBalancer}.

If peering fails, the controller returns a \texttt{RequeueAfter: 30s} result, causing the reconciliation loop to retry. Upon success, the instruction's \texttt{delivered} status is set to \texttt{true}.

\subsection{Result}

After successful peering, Liqo creates the following resources:
\begin{itemize}
    \item A \textbf{virtual node} in the requester cluster representing the provider's resources.
    \item A \textbf{WireGuard tunnel} for encrypted network communication.
    \item \textbf{Identity and authentication} resources for cross-cluster API access.
\end{itemize}

From this point, standard Kubernetes scheduling can place pods on the virtual node, and Liqo transparently offloads them to the provider cluster.


% =============================================================================
\section{Testing}
\label{sec:testing}

The implementation includes unit tests for the critical components:

\begin{itemize}
    \item \textbf{Decision engine:} Tests cluster selection with various resource distributions, edge cases (single cluster, all clusters full, self-exclusion), and priority handling. Uses mock \texttt{ClusterAdvertisement} resources with fake Kubernetes clients.

    \item \textbf{Resource calculation:} Tests the \texttt{CalculateAvailable} function with zero, positive, and negative (clamped) scenarios, including the Reserved field.

    \item \textbf{Authentication middleware:} Tests certificate CN extraction, missing certificate handling, health endpoint bypass, and cluster ID context propagation.

    \item \textbf{Metrics collector:} Tests node enumeration, pod request aggregation (including init container handling), and GPU resource tracking using fake Kubernetes clients.
\end{itemize}

Integration testing is performed through the test-setup scripts described in the \texttt{test-setup/} directory, which create three Kind clusters (one broker, two agents) and verify the complete reservation flow end-to-end. The evaluation suite (Chapter~\ref{ch:evaluation}) provides comprehensive system-level testing.
